%
% Informazioni us Sparse Autoencoders:
%   teoria degli autoencoder,
%   sparsity,
%   disentanglement,
%   mechanistic interpretability,
%   sparse features,
%   latent neurons.
%
\chapter{Sparse Autoencoders and Mechanistic Interpretability}

\section{Introduzione}

Uno dei principali problemi dei moderni Vision-Language Models riguarda la scarsa interpretabilità delle rappresentazioni latenti interne.

Come discusso nei capitoli precedenti, i neuroni presenti nei livelli nascosti delle reti neurali tendono a codificare simultaneamente molteplici concetti differenti, rendendo estremamente difficile comprendere il reale significato delle attivazioni interne.

Per affrontare questo problema, la ricerca recente si è concentrata su tecniche di:

\begin{itemize}
    \item \textit{Mechanistic Interpretability};
    \item \textit{Sparse Representation Learning};
    \item \textit{Disentanglement} delle rappresentazioni neurali.
\end{itemize}

In questo contesto, gli \textit{Sparse Autoencoders} (SAE) rappresentano uno degli approcci più promettenti per estrarre concetti interpretabili dai latent spaces dei Vision-Language Models.

\section{Teoria degli Autoencoder}

Prima di introdurre gli Sparse Autoencoders, è necessario comprendere il funzionamento di un Autoencoder standard.

Un \textit{Autoencoder} è una rete neurale progettata per apprendere una rappresentazione compatta dei dati di input.

L'obiettivo principale consiste nel ricostruire l'input originale dopo averlo fatto passare attraverso una rappresentazione compressa.

\section{Architettura di un Autoencoder}

Un Autoencoder è generalmente composto da due componenti principali:

\begin{itemize}
    \item un \textbf{Encoder};
    \item un \textbf{Decoder}.
\end{itemize}

\subsection{Encoder}

L'Encoder prende in input i dati originali e li trasforma in una rappresentazione latente compressa.

Formalmente:

\[
h = f(x)
\]

dove:

\begin{itemize}
    \item $x$ rappresenta l'input originale;
    \item $h$ rappresenta la rappresentazione latente;
    \item $f(\cdot)$ rappresenta la funzione appresa dall'encoder.
\end{itemize}

\subsection{Decoder}

Il Decoder riceve la rappresentazione latente e tenta di ricostruire l'input iniziale:

\[
r = g(h)
\]

dove:

\begin{itemize}
    \item $h$ è il vettore latente;
    \item $r$ rappresenta la ricostruzione finale;
    \item $g(\cdot)$ rappresenta la funzione del decoder.
\end{itemize}

L'obiettivo dell'addestramento consiste nel minimizzare la differenza tra input originale e ricostruzione:

\[
r \approx x
\]

\subsection{Bottleneck}

L'elemento fondamentale di un Autoencoder è il cosiddetto \textit{bottleneck}.

Poiché la rappresentazione latente ha dimensionalità ridotta rispetto all'input originale, la rete non può semplicemente memorizzare i dati.

Il modello è quindi costretto a:

\begin{itemize}
    \item apprendere pattern rilevanti;
    \item identificare strutture statistiche significative;
    \item comprimere le informazioni più importanti.
\end{itemize}

\section{Mechanistic Interpretability}

La \textit{Mechanistic Interpretability} rappresenta un ramo recente dell'Explainable AI che cerca di analizzare le reti neurali come sistemi matematici interpretabili.

L'obiettivo consiste nel trattare le reti neurali come sistemi analizzabili tramite \textit{reverse engineering}, analogamente a:

\begin{itemize}
    \item programmi software compilati;
    \item circuiti elettronici;
    \item sistemi biologici complessi.
\end{itemize}

Si cerca di comprendere:

\begin{itemize}
    \item il ruolo dei singoli neuroni;
    \item le trasformazioni effettuate dai layer;
    \item la struttura interna delle rappresentazioni latenti.
\end{itemize}

\subsection{Neuroni Latenti}

Nei Vision-Language Models, i neuroni presenti nei livelli nascosti costituiscono le cosiddette rappresentazioni latenti.

Tali neuroni codificano informazioni astratte relative a:

\begin{itemize}
    \item forme;
    \item texture;
    \item pattern anatomici;
    \item concetti semantici;
    \item relazioni statistiche.
\end{itemize}

\subsection{Polisemanticità}

Uno dei principali problemi riguarda la \textit{polisemanticità} dei neuroni.

Un singolo neurone può infatti attivarsi simultaneamente rispetto a concetti completamente differenti.

Ad esempio, un neurone potrebbe rispondere contemporaneamente a:

\begin{itemize}
    \item presenza di una frattura;
    \item specifiche texture visive;
    \item luminosità elevate;
    \item forme anatomiche.
\end{itemize}

Questa sovrapposizione semantica rende estremamente difficile interpretare il reale significato delle attivazioni interne.

\section{Sparse Autoencoders}

Per affrontare il problema della polisemanticità, la ricerca recente ha introdotto gli \textit{Sparse Autoencoders}.

L'obiettivo principale dei SAE consiste nel:

\begin{itemize}
    \item separare concetti differenti;
    \item creare rappresentazioni sparse;
    \item ottenere feature maggiormente interpretabili.
\end{itemize}

\subsection{Differenze rispetto agli Autoencoder Tradizionali}

A differenza degli Autoencoder classici, gli Sparse Autoencoders non comprimono necessariamente l'informazione in uno spazio più piccolo.

Molti SAE utilizzano invece un approccio chiamato:

\begin{quote}
\textit{Overcomplete Dictionary Learning}
\end{quote}

In questo scenario:

\begin{itemize}
    \item un embedding originale può avere dimensione 512;
    \item il SAE può espandere tale rappresentazione a migliaia di dimensioni;
    \item ogni dimensione può rappresentare una feature semantica specifica.
\end{itemize}

Ad esempio:

\[
512 \rightarrow 16000
\]

L'espansione dello spazio latente fornisce al modello sufficiente capacità per separare concetti differenti.

\section{Sparsity}

L'espansione dello spazio latente, da sola, non è sufficiente per ottenere interpretabilità.

Per questo motivo viene introdotto un vincolo fondamentale:

\begin{quote}
\textbf{la sparsità}
\end{quote}

\subsection{Vincolo di Sparsità}

Durante il training del SAE viene aggiunta una penalizzazione matematica, tipicamente basata sulla norma $L_1$:

\[
\lambda \|h\|_1
\]

Questo termine forza il modello a utilizzare solo un numero molto limitato di neuroni attivi per ogni input.

L'obiettivo è ottenere rappresentazioni in cui:

\begin{itemize}
    \item pochi neuroni risultano attivi;
    \item la maggior parte delle feature rimane pari a zero.
\end{itemize}

La penalizzazione forza il modello a:

\begin{itemize}
    \item utilizzare pochissime feature contemporaneamente;
    \item mantenere inattiva la maggior parte dei neuroni;
    \item produrre attivazioni sparse.
\end{itemize}

\subsection{Sparse Features}

Le feature generate dagli Sparse Autoencoders vengono chiamate \textit{Sparse Features}.

L'idea fondamentale è che ogni feature possa rappresentare un singolo concetto semanticamente coerente.

Ad esempio:

\begin{itemize}
    \item una feature può rappresentare un versamento pleurico;
    \item un'altra può rappresentare il polmone destro;
    \item un'altra ancora può codificare specifiche anomalie anatomiche.
\end{itemize}

In teoria, le sparse features risultano quindi molto più interpretabili rispetto ai neuroni polisemantici originali.

\section{Disentanglement}

Il risultato finale del processo di sparsificazione prende il nome di \textit{Disentanglement}.

Il termine indica la capacità di:

\begin{itemize}
    \item separare concetti mescolati;
    \item isolare fattori latenti indipendenti;
    \item ottenere rappresentazioni semanticamente organizzate.
\end{itemize}

Nei Vision-Language Models originali, differenti concetti risultano fortemente intrecciati all'interno degli stessi neuroni.

Gli Sparse Autoencoders cercano invece di:

\begin{itemize}
    \item districare tali rappresentazioni;
    \item associare concetti differenti a feature differenti;
    \item produrre attivazioni maggiormente monosemantiche.
\end{itemize}

\section{Collegamento con MedConcept}

Il framework \textit{MedConcept} utilizza proprio Sparse Autoencoders per estrarre concetti latenti interpretabili dai Vision-Language Models medici.

L'obiettivo è ottenere:

\begin{itemize}
    \item scoperta non supervisionata di concetti;
    \item interpretabilità neuronale;
    \item allineamento semantico con terminologia clinica;
    \item maggiore trasparenza delle rappresentazioni latenti.
\end{itemize}

A differenza di approcci supervisionati come TCAV o Concept Bottleneck Models, MedConcept non richiede annotazioni manuali preventive.

I concetti vengono invece scoperti automaticamente direttamente dalle attivazioni neurali del modello.

Questo rappresenta uno degli aspetti più innovativi e promettenti dell'intero framework.